\section{The Age of the Feminine}

Recently, someone objected that Gornahoor overemphasizes hierarchy, authority, and intelligence at the moment that ``we are entering an age, where feminine values of the heart, a sense of human equality and individuality are becoming ever more predominant.''

Since the comments were made in an intelligent and civil manner, I want to explore that idea in a sustained way. Clearly, we both agree that the present age emphasizes ``feminine'' values, egalitarianism, and individualism; however, Gornahoor regards that as a symptom of decline, unlike the majority who are caught up in the spirit of the age. Unfortunately, a polemical attitude tends to make a man see only the worst in the position he opposes, and the best in what he proposes. The Hermetic method, on the other hand, seeks depth, that is, to find the common principle. So, instead of opposing the Masculine to the Feminine, we try to understand the principle of unity, or the Tao. Within this principle, we will find the correct relationship between Masculine and Feminine ``values''.

Our interlocutor objects that he sees nothing to justify a hierarchy at the present time. As far as it goes, that is correct; but that is not because our present age is evolving, but rather just the opposite. The modern mind objects that ``blood does not make the noble''. Again correct, but the man of Tradition says that a man will act nobly because of his blood. The modern mind objects that today's spiritual leaders are not such simply by virtue of their station. Of course, the Medieval tradition based the spiritual hierarchy on the order of merit, which would be open to anyone so qualified.

The difference, then, in our perspectives is that we do not blindly resign ourselves to current conditions; rather we look to restore Tradition. This is not at all the same as nostalgia for a lost time. Rather, if the pretenders to spiritual authority and temporal power lack the requisite qualification, then they must be replaced. This does not mean at all that authority and power are defective as concepts---an absurd idea since they will reappear in unexpected ways. It means that the men holding those positions need to be replaced. If the masculine age has given way to the feminine, as he claims, then it is just as possible for the masculine age of hierarchy and authority to also return.

We can immediately question the idea of equality, which is only possible in the Primordial State. \textbf{Ezekiel} spoke of the righteousness of \textbf{Noah}, \textbf{Job}, and \textbf{Daniel}. (In passing, we can mention the importance this prophet had for \textbf{Valentin Tomberg}.) This idea came to prominence in the Middle Ages from Gregory's commentary on that book; hence, rank in the spiritual hierarchy had to be based on merit. This was not a bane. In the chaos of the fall of Rome, order was established based on merit. Those who fight, and not all are capable of that, especially women, were necessary to enforce order out of that chaos. The spiritual authority guaranteed that the rights of the poor, the widows, and the orphans were protected.

Could that have been accomplished on the basis of the values of fraternity, equality, and freedom? The Terror of the French Revolution was necessary after the overthrow of all authority. Tomberg considered that the French revolution was an ``orgy of perverse collective intoxication'' and the result of a ``dark current''. Since the events of 1789 form the philosophical basis of the modern world, that is exactly how a well-bred man of today must evaluate the ``predominant'' features of this world.

As for the ``values of the heart'', this is based on a misunderstanding. The heart refers to the central core of man, his Spirit or Intellectual Soul, and does not mean sentimentalism as the modern mind believes. \textbf{Frithjof Schuon} explains the reason for this common error:


\begin{quotex}
The use, in the esoterisms, of the term ``love'' to designate an intellectual reality is moreover explained by the fact that sentiment, while being inferior to reason because of its emotional subjectivity, is nevertheless symbolically comparable to what is superior to reason, namely the Intellect; it is so because sentiment, like the Intellect at whose antipodes it stands, is not discursive, but direct, simple, spontaneous, unlimited, compared with reason, sentiment appears free from form and fallibility, and this is why the Divine Intelligence can be called ``Love''. Modes of Spiritual Realization \end{quotex}


However, sentiment is passive and, again quoting Schuon:

\begin{quotex}
Only he who possesses the truth in an active manner is really intellectual, and not he who accepts it passively ... \end{quotex}


The modern world has lost its intellectuality and cannot understand itself. Instead, every idea is related to ``feelings'', which come, go, and change. This passivity of sentiment may seem on the surface to be caring and kind, but it can turn vicious against those who see past the appearances.

The feminine is the reflection of the masculine, it is an illusion to be on its own without the masculine element, or worse, to be in opposition. Tomberg says that the masculine is Love, since God is Love (see Schuon quote above), and its reflection, as the feminine aspect, is Wisdom. This is the opposite of what is commonly believed, but, once again, that is due to an overly sentimental understanding of Love.

Love is the unifying force, and only the Intellect can see individual things in their wholeness. Hence, a world of ``individuals'' makes no sense. To believe in individuals is to believe that a collection of sticks makes a house. No, only sticks that are related to each other in a specific way can create a home. The Person, on the other hand, is related: to his family, his nation, his civilization. These are hierarchical relations, quite different from the egalitarian relationships of sticks.

The Traditional way is to see society as an organic whole, with functions and roles working harmoniously together. Harmony is concord, whose root meaning is ``one heart''; thus an organic society is based on one heart, one mind, one understanding. This is not mindless conformance, but is organized so that each member can best fit its role. Its opposite is the feminine mass culture of today that forces politically correct conformance and uses the beehive or the anthill as its model of perfectly functioning individuals.

To believe in the Age of the Feminine, is to believe in the reflection without believing in its source; it is to live in Plato's cave believing the shadows to be real and not knowing the real things that cast them. Shadows are necessarily two dimensional. The most fantastic skyline still casts a flat shadow. To believe in the shadows is to miss the third dimension, the hierarchy that reaches up to the heavens. That is our situation today.

This is not to reject the Feminine, because as Divine Sophia, She offers the gift of \emph{sapientia}, a gaze capable of penetrating behind the veil of appearances to reach hidden truths. \textbf{Blessed Adalbero}\footnote{\url{https://www.catholic.org/encyclopedia/view.php?id=393}}


\flrightit{Posted on 2012-11-20 by Cologero}

\begin{center}* * *\end{center}

\begin{footnotesize}
\begin{sffamily}
\texttt{Logres on 2012-11-21 at 06:26 said:}

Samael Aun Weor remarks in Perfect Matrimony that a true marriage is between the one who ``loves most'' (male -- Love) and the one ``who loves better'' (women -- wisdom).

\hfill

\texttt{Maude Murray on 2022-11-20 at 05:00 said:}

I feel as if I've, by God's will, seen Purgatory and He'll, but am now ``with Beatrice.'' Life has been full of tragedies and disappointment in people I thought were saints, and then refused to believe otherwise until dead! For others, because life in these times is really hard for spiritual people, I want to share something that has ALWAYS PROVEN TRUE: ``We know that for those who love God, all things work together unto good... (Romans..).Every price of a pin, or blow of a sword, has become good, brought wisdom, good people, deeper faith!

\hfill

\texttt{Arthur Konrad on 2022-11-20 at 14:12 said:}

Mrs. Murray,

Fate was very kind to you if it brought you to be in the company such illustrious people and sages. This perhaps applies to men -- unless one belongs to those few born with a god-like steadfastness, misfortune is essentially a mystery one must brace in order to know anything about himself. Refusal to deal with it invokes only more misery, cause there is no advancing without having learnt the lesson of today. No reason to resign oneself to living in the shades!

Finest things are a secret. Muslims understand this well. The value of piety is nothing. Loving in secret is a fountain of good fortune.


\hfill

\end{sffamily}
\end{footnotesize}

